We define AI hallucinations as a manifestation of \textit{Proofreading Deficiency}. In biological systems, the inability of DNA Polymerase to excise mismatched bases leads to genomic instability. Similarly, in LLMs, the lack of an intrinsic self-audit mechanism during token generation leads to logical divergence.

\subsection{Digital Pathogenesis}
We treat reasoning drift as a systemic disease. A ``hallucination'' is not an isolated error but a symptom of a failing context-identity. By adopting an \textbf{organism-mimetic approach}, we can apply adaptive immunity to recognize ``non-self'' logic (pathogenic code) and eliminate it through dialectical synthesis.

\subsection{Mechanistic Divergence: Immune vs. Mutation}
We argue that a hallucination is ontologically dual-natured:
\begin{itemize}
    \item \textbf{The Genetic Mutation Model (Micro)}: Focuses on sequence fidelity at the point of synthesis (Vibe-to-Code translation).
    \item \textbf{The Immune Response Model (Macro)}: Focuses on systemic integrity and logical alignment across the entire reasoning chain.
\end{itemize}
The convergence of these two models ensures both local stability and global alignment.
